% TEMPLATE-MILC-v3.tex - plantilla maestra UNICA de las guias MILC v3.
%
% La usan las cuatro familias de guias, cada una con su driver:
%   · Grados 6-11          -> scripts/build-guias-g11.py
%   · Bebras               -> scripts/build-guias-bebras.py
%   · Semillero            -> scripts/build-guias-semillero.py
%   · Territorio interior  -> scripts/build-guias-territorio-interior.py
%
% Antes habia tres copias casi identicas (~400 lineas iguales, 6 distintas):
% cada mejora habia que aplicarla tres veces, y en la practica no se hacia
% (el motor de recursos vivia solo en dos de ellas). Lo que varia entre
% familias esta parametrizado en 6 placeholders que cada driver llena
% (se nombran aqui SIN su sintaxis real: el builder reemplaza por texto plano,
% tambien dentro de los comentarios, y un valor multilinea rompe el preambulo):
%
%   HEADER_IZQ        encabezado izquierdo de pagina
%   HEADER_DER        encabezado derecho de pagina
%   PORTADA_ETIQUETA  rotulo sobre el numero grande (GRADO/MOMENTO/MODULO)
%   PORTADA_SERIE     linea de serie de la portada
%   PORTADA_ID        identificador de la guia en la portada
%   PORTADA_CIRCULO   el circulito de grado (°), o vacio si no aplica
%   PDF_TITULO        titulo en texto plano (sin \textbf ni comillas TeX) para
%                     los metadatos del PDF (/Title); lo produce lib_xelatex.texto_pdf
%
% Composicion (auditoria 2026-09-03): las bandas de estacion piden espacio
% (\Needspace*) y prohiben el salto justo despues (\nopagebreak), asi no quedan
% huerfanas al pie; las cajas largas (softbox, drawbox, infoband, puentebox) son
% `breakable`, asi una caja de media pagina ya no arrastra todo a la siguiente
% dejando la anterior al 40 %. Todo texto sobre mostaza o verde va en milcNegro
% (blanco sobre #E5B400 da 1,93:1 y sobre #58A923 2,95:1; ninguno pasa WCAG AA).
% El resto de claves las documenta content/guias/_SCHEMA.md. El script valida
% que no queden placeholders sin reemplazar antes de escribir el archivo final.

% Etiquetado PDF (latex-lab): /Lang, estructura y texto alternativo de las
% figuras, para lectores de pantalla. Debe ir ANTES de \documentclass.
\DocumentMetadata{lang=es-CO,tagging=on}
\documentclass[11pt,letterpaper]{article}
% headheight=24pt: la cabecera derecha lleva dos lineas (identidad de la guia
% y estacion actual). headsep se acorta en la misma medida para que el bloque
% de texto quede exactamente donde estaba con headheight=12pt.
\usepackage[margin=2.0cm,headheight=24pt,headsep=13pt]{geometry}
\usepackage[table]{xcolor}
\usepackage{fontspec}
\usepackage[spanish]{babel}
\usepackage{tikz}
% Librerías TikZ para los diagramas de las guías (bloque `recursos:`):
%  · babel  → babel en español vuelve activos caracteres como «>», y TikZ los usa
%             en sus opciones (>=stealth, ->). Sin esta librería, cualquier
%             diagrama con flechas revienta con «\language@active@arg> has an
%             extra }». Debe ir ANTES de que se dibuje nada.
%  · positioning        → sintaxis «below left=2pt and 3pt».
%  · decorations.pathreplacing → llaves ({brace}) para señalar rangos.
\usetikzlibrary{babel,positioning,decorations.pathreplacing}
\usepackage{graphicx}
% Circuitos: el trabajo de electrónica se hace hoy con micro:bit y sus sensores
% integrados (MakeCode, sin cableado), así que no se necesitan esquemáticos.
% Si algún día entran montajes con protoboard, basta descomentar esta línea:
% los diagramas .tex/.tikz declarados en `recursos:` ya se incrustan solos.
% \usepackage{circuitikz}
\usepackage[most]{tcolorbox}
\usepackage{tabularx}
\usepackage{array}
\usepackage{fancyhdr}
\usepackage{titlesec}
\usepackage[document]{ragged2e}  % bandera a la derecha en todo el cuerpo
\usepackage{enumitem}
\usepackage{microtype}
\usepackage{etoolbox}
\usepackage{needspace}
% hyperref al final del preambulo: metadatos del PDF (titulo, autor, idioma,
% familia), marcadores por estacion (\pdfbookmark) y enlaces sin recuadro.
\usepackage[hidelinks,
  pdftitle={Proyecto de monitoreo — tres jornadas de datos},
  pdfauthor={PhD. Álvaro Cárdenas Orozco},
  pdfsubject={Tecnología e Informática · Grado 8 · Guía 19 · MILC v3},
  pdflang={es-CO},
  bookmarksopen=true,bookmarksnumbered=false]{hyperref}

% Helvetica Neue no trae la flecha U+2192, asi que cada «→» escrito literalmente
% en el YAML se compilaba a NADA, en silencio (462 flechas en 61 guias: rutas del
% tipo «paleta → tipografia → lockup» salian rotas). Mapeamos el caracter a su
% version matematica, que si tiene glifo. Asi el autor puede seguir escribiendo
% «→» comodamente en el YAML.
\usepackage{newunicodechar}
\newunicodechar{→}{\ensuremath{\rightarrow}}
% Mismo problema con los simbolos de marca y vinetas: el «✓» de «una palomita ✓»
% salia como cuadrito vacio en el PDF de todas las guias que lo usaban.
\usepackage{pifont}
\newunicodechar{✓}{\ding{51}}
\newunicodechar{✗}{\ding{55}}
\newunicodechar{✔}{\ding{52}}
\newunicodechar{•}{\textbullet}
\newunicodechar{≥}{\ensuremath{\geq}}
\newunicodechar{≤}{\ensuremath{\leq}}

\setmainfont{Helvetica Neue}
\setsansfont{Helvetica Neue}
\setlength{\parindent}{0pt}
\setlength{\parskip}{6pt}
% Sin particion de palabras: en 6.º-7.º una palabra cortada por guion al final
% de linea («Comuni-cacion») frena la lectura mucho mas que una linea corta.
\hyphenpenalty=10000
\exhyphenpenalty=10000
\pretolerance=2000
\tolerance=3000
\setlength{\emergencystretch}{3em}
\linespread{1.10}
\renewcommand{\arraystretch}{1.25}
\setlist{leftmargin=5mm,itemsep=1mm,topsep=1mm}
\newcolumntype{Y}{>{\RaggedRight\arraybackslash}X}

\definecolor{milcMagenta}{HTML}{E6007E}
\definecolor{milcVino}{HTML}{5A0038}
\definecolor{milcTurquesa}{HTML}{0093A5}
\definecolor{milcVerde}{HTML}{58A923}
\definecolor{milcMostaza}{HTML}{E5B400}
\definecolor{milcOcre}{HTML}{B86600}
\definecolor{milcNegro}{HTML}{241816}
\definecolor{milcGris}{HTML}{F7F7F4}
\definecolor{milcLinea}{HTML}{E8E2DD}
\definecolor{milcGradePrimary}{HTML}{FF6600}
\definecolor{milcGradeDark}{HTML}{9A3E00}
\definecolor{milcGradeSoft}{HTML}{FFE4D0}
\definecolor{milcGradeAccent}{HTML}{CFE7FF}
\definecolor{milcGradeText}{HTML}{FFFFFF}
\color{milcNegro}

\pagestyle{fancy}
\fancyhf{}
% Cabecera de dos lineas: identidad de la guia arriba y, a la derecha y en
% gris, la estacion en curso (\markright desde \milcestacion). Antes de la
% primera estacion la segunda linea va vacia; el \strut mantiene la altura
% para que la linea de arriba no baile entre paginas.
\lhead{\small Tecnología e Informática · Grado 8\\[-1pt]{\footnotesize\strut}}
\rhead{\small Guía 19 · MILC v3\\[-1pt]{\footnotesize\strut\color{milcNegro!65}\rightmark}}
\cfoot{\small\thepage}
\renewcommand{\headrulewidth}{0pt}

% Sobre mostaza (#E5B400) y verde (#58A923) el texto va en milcNegro; sobre el
% resto de la paleta, en blanco. Se usa dentro de fonttitle/fontupper, donde un
% \color posterior manda sobre coltitle.
\newcommand{\milctintasobre}[1]{%
  \ifboolexpr{test{\ifstrequal{#1}{milcMostaza}} or test{\ifstrequal{#1}{milcVerde}}}%
    {\color{milcNegro}}{\color{white}}}

\newtcolorbox{titlebox}[1]{
  enhanced,colback=#1,colframe=#1,arc=7mm,boxrule=0pt,
  left=7mm,right=7mm,top=3mm,bottom=3mm,
  before skip=8pt,after skip=8pt,
  fontupper=\sffamily\bfseries\Large\color{white}
}

\newtcolorbox{softbox}[3]{
  enhanced,breakable,pad at break=3mm,
  colback=#2,colframe=#1,borderline west={4pt}{0pt}{#1},
  arc=2mm,boxrule=.4pt,left=4mm,right=4mm,top=3mm,bottom=3mm,
  before skip=6pt,after skip=8pt,title={#3},coltitle=white,
  colbacktitle=#1,fonttitle=\sffamily\bfseries\milctintasobre{#1},fontupper=\RaggedRight,
  attach boxed title to top left={xshift=3mm,yshift=-2mm},
  boxed title style={arc=2mm, boxrule=0pt}
}

\newtcolorbox{drawbox}[2]{
  enhanced,breakable,pad at break=4mm,
  colback=white,colframe=#1,arc=4mm,boxrule=1pt,
  left=5mm,right=5mm,top=4mm,bottom=4mm,before skip=8pt,after skip=8pt,
  title={#2},coltitle=white,colbacktitle=#1,fonttitle=\sffamily\bfseries\milctintasobre{#1},
  fontupper=\RaggedRight,
  attach boxed title to top left={xshift=4mm,yshift=-2mm},
  boxed title style={arc=3mm, boxrule=0pt}
}

\newtcolorbox{infoband}[1]{
  enhanced,breakable,pad at break=2mm,colback=#1!8,colframe=#1!50,arc=2mm,boxrule=.5pt,
  left=4mm,right=4mm,top=2mm,bottom=2mm,before skip=4pt,after skip=4pt,
  fontupper=\small\RaggedRight
}

% Puente narrativo entre secciones (contrato editorial Conéctate).
% Da continuidad: "Acabas de X. Ahora vas a Y porque Z."
% Visualmente: banda izquierda en color del grado, texto en itálica pequeña.
\newtcolorbox{puentebox}{
  enhanced,breakable,pad at break=2mm,colback=milcGradeSoft,colframe=milcGradeDark,
  borderline west={3pt}{0pt}{milcGradeDark},
  arc=0mm,boxrule=0pt,left=5mm,right=4mm,top=2.5mm,bottom=2.5mm,
  before skip=4pt,after skip=8pt,
  fontupper=\itshape\small\color{milcNegro}\RaggedRight
}
% La flecha va en modo matematico a proposito: Helvetica Neue NO tiene el glifo
% U+2192, asi que \textrightarrow se compilaba a NADA («Missing character: There
% is no → in font Helvetica Neue») y el puente salia sin flecha. En modo
% matematico la flecha viene de la fuente matematica y si se dibuja.
\newcommand{\puente}[1]{%
  \begin{puentebox}%
    \textcolor{milcGradeDark}{$\rightarrow$}\hspace{2mm}#1%
  \end{puentebox}%
}

\newcommand{\linead}[1]{\noindent\color{milcLinea}\rule{#1}{0.5pt}\color{milcNegro}}
\newcommand{\respuesta}[1]{\par\vspace{2mm}\foreach \n in {1,...,#1}{\noindent\color{milcLinea}\rule{\linewidth}{0.5pt}\par\vspace{5mm}}\color{milcNegro}}
\newcommand{\checkbox}{\textcolor{milcGradeDark}{\fbox{\phantom{x}}}\hspace{5pt}}

% ─── Listas aireadas para las guías socioemocionales ────────────────────
% Componer las verificaciones como «\checkbox texto\\» deja el texto que salta
% de línea debajo del cuadrito y apelmaza el bloque. Estos entornos alinean la
% continuación y airean los ítems. Los usa Territorio Interior; cualquier
% familia puede hacerlo.
\newlist{checklista}{itemize}{1}
\setlist[checklista]{
  label=\textcolor{milcGradeDark}{\fbox{\phantom{x}}},
  leftmargin=9mm, labelsep=3.2mm, itemsep=2.4mm,
  topsep=2.5mm, parsep=0pt, partopsep=0pt, before=\RaggedRight
}
\newlist{pasoslista}{enumerate}{1}
\setlist[pasoslista]{
  label=\textcolor{milcGradeDark}{\sffamily\bfseries\arabic*.},
  leftmargin=8mm, labelsep=3mm, itemsep=2.8mm,
  topsep=1mm, parsep=0pt, partopsep=0pt, before=\RaggedRight
}

\newcommand{\phasecard}[4]{
\begin{tcolorbox}[enhanced,colback=#3,colframe=#2,arc=3mm,boxrule=.5pt,height=3.05cm,valign=center,halign=center,left=1.5mm,right=1.5mm,top=2mm,bottom=2mm]
{\sffamily\bfseries\fontsize{22}{24}\selectfont\textcolor{#2}{#1}}\par
{\sffamily\bfseries\small #4}
\end{tcolorbox}
}

% ── Piezas del rediseno 2026 (legibilidad 6.º-11.º) ───────────────────
% Banda de estacion: reemplaza el titlebox macizo. Titulo a la izquierda,
% dato practico (tiempo/modalidad) a la derecha, en una sola linea.
% Nunca huerfana: pide 9 lineas libres (o salta de pagina rellenando con
% \vfill) y prohibe el salto justo despues de la banda. Nueve y no seis:
% la caja partible que sigue necesita ~80pt (titulo adosado, relleno y dos
% lineas) para arrancar en la misma pagina; con menos, tcolorbox la manda
% entera a la siguiente y la banda se queda sola. El penalty 10000 va
% en `after`, ANTES del espacio posterior: un `after skip` (aunque sea 0pt)
% es un \vskip tras la caja, y ese glue es un punto de corte legal que un
% \nopagebreak escrito despues de \end{tcolorbox} ya no cierra. Cada banda
% deja marcador PDF y actualiza la cabecera derecha.
\newcounter{milcestacion}
\newcommand{\milcestacion}[2]{%
\Needspace*{9\baselineskip}%
\stepcounter{milcestacion}%
\pdfbookmark[1]{#2}{milcestacion\themilcestacion}%
\markright{#2}%
\begin{tcolorbox}[enhanced,colback=#1,colframe=#1,arc=2mm,boxrule=0pt,
  left=5mm,right=5mm,top=2.6mm,bottom=2.6mm,before skip=11pt,
  after={\par\nopagebreak\vskip7pt}]
{\sffamily\bfseries\fontsize{15}{18}\selectfont\milctintasobre{#1}#2}
\end{tcolorbox}}

% Tarjeta de paso numerada: sustituye las tablas «Pilar / Aplicacion», que
% eran formato de consulta docente, no de estudiante. El numero va en un
% circulo sobre el borde para que la secuencia se lea de un vistazo.
\newcommand{\milcpaso}[3]{%
\Needspace{4\baselineskip}%
\begin{tcolorbox}[enhanced,colback=white,colframe=milcLinea,arc=1.5mm,boxrule=.9pt,
  left=14mm,right=4mm,top=3mm,bottom=3mm,before skip=4pt,after skip=4pt,
  fontupper=\RaggedRight,
  overlay={\node[anchor=north west,circle,fill=#2,text=white,inner sep=0pt,
    minimum size=7.4mm,font=\sffamily\bfseries\small]
    at ([xshift=3.6mm,yshift=-3.2mm]frame.north west) {#1};}]
#3
\end{tcolorbox}}

% Casilla marcable de tamano util con lapiz (la anterior era \fbox{\phantom{x}},
% de alto variable segun el texto que la seguia).
\newcommand{\milcmarca}{\framebox[3.8mm]{\rule{0pt}{3.4mm}}}

% Fila marcable de lista suelta (checks de la estacion 1).
\newcommand{\milccheckrow}[1]{%
\par\vspace{1.8mm}\noindent
\makebox[6.4mm][l]{\raisebox{-.55mm}{\textcolor{milcNegro}{\milcmarca}}}%
\parbox[t]{\dimexpr\linewidth-6.4mm\relax}{\RaggedRight #1}\par}

% Celda marcable dentro de la rubrica (tabla de autoevaluacion). Misma
% sangria francesa, pero pensada para vivir dentro de una columna X.
\newcommand{\milcopcion}[1]{%
\makebox[5.2mm][l]{\raisebox{-.55mm}{\textcolor{milcNegro}{\milcmarca}}}%
\parbox[t]{\dimexpr\linewidth-5.2mm\relax}{\RaggedRight #1}}

% ── Recursos visuales de la guía ──────────────────────────────────────
% Declarados en el bloque `recursos:` del YAML y emitidos por el builder.
% \guiaFigura   → imagen rasterizada o PDF (foto, carta celeste, montaje).
% \guiaDiagrama → fuente TikZ/circuitikz (.tex) incrustada como vector:
%                 es la vía para circuitos y esquemáticos editables.
% [#1] = texto alternativo (alt) · #2 = ruta relativa al .tex · #3 = pie
% El alt llega al PDF etiquetado (/Alt) y lo lee el lector de pantalla.
\newcommand{\guiaFigura}[3][]{%
\begin{center}
\includegraphics[width=0.88\linewidth,height=0.38\textheight,keepaspectratio,alt={#1}]{#2}\\[1.5mm]
{\footnotesize\itshape\textcolor{milcNegro!75}{#3}}
\end{center}
}

\newcommand{\guiaDiagrama}[2]{%
\begin{center}
\input{#1}\\[1.5mm]
{\footnotesize\itshape\textcolor{milcNegro!75}{#2}}
\end{center}
}

\begin{document}
\thispagestyle{empty}

% ── Portada ───────────────────────────────────────────────────────────
% Antes: un rectangulo de color a pagina completa. Se veia bien en pantalla,
% pero en la fotocopiadora del colegio se comia el toner y salia gris sucio.
% Ahora la banda ocupa el tercio superior y el resto va en blanco.
\begin{tikzpicture}[remember picture,overlay]
  \path[fill=milcGradePrimary]
    (current page.north west) rectangle ([yshift=-10.6cm]current page.north east);

  \node[anchor=north west,text=milcGradeText] at ([xshift=2.0cm,yshift=-1.55cm]current page.north west)
    {\sffamily\bfseries\fontsize{12}{14}\selectfont GRADO};
  \node[anchor=north west,text=milcGradeText,opacity=.85] at ([xshift=2.0cm,yshift=-2.35cm]current page.north west)
    {\sffamily\bfseries\fontsize{8.5}{10}\selectfont SERIE GUÍAS · TECNOLOGÍA E INFORMÁTICA};
  \node[anchor=north east,text=milcGradeText,opacity=.85,align=right] at ([xshift=-2.0cm,yshift=-1.55cm]current page.north east)
    {\sffamily\bfseries\fontsize{9}{12}\selectfont I.E. Sor María Juliana\\Cartago, Valle del Cauca};

  \node[anchor=west,text=milcGradeText] (gradonum) at ([xshift=1.85cm,yshift=-7.1cm]current page.north west)
    {\sffamily\bfseries\fontsize{112}{112}\selectfont 8};
  % El circulito de grado (°) solo aplica a las guias de grado («11°»); en
  % Bebras y Semillero el numero grande es un momento/modulo, no un grado.
  % Se posiciona relativo al nodo `gradonum`, no en coordenadas absolutas.
  \draw[milcGradeText,line width=1.6pt]
    ([xshift=2.0mm,yshift=-4.5mm]gradonum.north east) circle (.15cm);

  \node[anchor=west,text=milcGradeText,text width=11.4cm,align=left] at ([xshift=1.5cm,yshift=0.5cm]gradonum.east)
    {
     {\sffamily\bfseries\fontsize{9}{11}\selectfont\MakeUppercase{Período 2 · Lógica, algoritmos y computación física}}\\[3mm]
     {\sffamily\bfseries\fontsize{21}{25}\selectfont\setlength{\baselineskip}{27pt}Tres jornadas\\[4pt]de datos}\\[3.5mm]
     {\sffamily\bfseries\fontsize{15}{18}\selectfont Guía 19-8-TIC}
    };
\end{tikzpicture}

\vspace*{9.4cm}
\vfill

{\sffamily\bfseries\fontsize{8}{10}\selectfont\textcolor{milcGradeDark}{LO QUE TE LLEVAS HOY}}

\begin{tcolorbox}[enhanced,colback=white,colframe=milcNegro,arc=2mm,boxrule=1.6pt,
  left=5mm,right=5mm,top=4mm,bottom=4mm,before skip=3pt,after skip=12pt,
  fontupper=\RaggedRight\fontsize{11}{15.5}\selectfont]
Un micro:bit programado para medir una variable real del colegio, temperatura del aula, luz del salón, ruido del recreo o paso de personas, durante tres jornadas. Una bitácora con al menos treinta mediciones, diez por jornada, con hora, valor, observación y nivel de alerta. Y un análisis de un patrón, la cifra que lo respalda y una propuesta de acción que salga de los datos.
\end{tcolorbox}

\begin{tcolorbox}[enhanced,colback=milcGris,colframe=milcGris,arc=2mm,boxrule=0pt,
  left=5mm,right=5mm,top=3.5mm,bottom=3.5mm,after skip=12pt]
{\sffamily\bfseries\fontsize{8}{10}\selectfont\textcolor{milcNegro!55}{ESTA GUÍA ES DE}}\\[3mm]
\begin{tabularx}{\linewidth}{X p{3.2cm} p{3.2cm}}
\linead{\linewidth} & \linead{3.0cm} & \linead{3.0cm}\\[-1mm]
{\fontsize{8.5}{10}\selectfont\textcolor{milcNegro!55}{Nombre completo}} &
{\fontsize{8.5}{10}\selectfont\textcolor{milcNegro!55}{Grupo}} &
{\fontsize{8.5}{10}\selectfont\textcolor{milcNegro!55}{Fecha}}\\
\end{tabularx}
\end{tcolorbox}

\vfill

\noindent\textcolor{milcLinea}{\rule{\linewidth}{1pt}}\\[2mm]
{\sffamily\bfseries\fontsize{9.5}{12}\selectfont Autor: Dr. Álvaro Cárdenas Orozco}\\[1mm]
{\sffamily\fontsize{8.5}{11}\selectfont\textcolor{milcNegro!55}{Metodología MILC · apertura ancestral · pensamiento computacional · triángulo Dussel–estoicismo–Floridi}}

\newpage

\section*{Apertura: Proyecto de monitoreo --- tres jornadas de datos}

\begin{softbox}{milcGradeDark}{milcGradeSoft}{Saber ancestral}
El café no se recoge todo el año igual. En Colombia hay dos cosechas: la principal, hacia el final del año, y una segunda más pequeña, la \textbf{mitaca} o traviesa, hacia mitad de año. El Valle del Cauca está en traviesa en el primer semestre mientras otras regiones recogen la grande. Entre una cosecha y otra, los campesinos de las veredas cafeteras de Caldas hablan del \textbf{tiempo frío}: los meses en que el cafetal no da fruto (Parada Sanabria, 2017). Ese tiempo no es tiempo perdido. Es la otra mitad del ciclo, la que hace posible la siguiente cosecha. Y solo se entiende mirando el año completo. Quien mira el cafetal una semana ve un cafetal sin grano; quien lo mira doce meses ve dos cosechas y un descanso. Un patrón necesita tiempo para aparecer. La \textbf{cara de exclusión}: el descanso es del cultivo, no del recolector; en tiempo frío no hay jornal, y quien vive de recoger café pasa esos meses buscando otro trabajo. Hoy vas a dejar el micro:bit midiendo durante tres jornadas, porque treinta datos en tres días muestran lo que cien datos en una hora esconden.\par\smallskip{\footnotesize\textcolor{milcNegro!70}{\textbf{Fuente.} Parada Sanabria, P. J. (2017). Práctica social y cultural del campesinado cafetero en cuatro municipios de Caldas (Colombia). \emph{Revista Colombiana de Sociología, 40}(1 Supl.), 193--212. https://doi.org/10.15446/rcs.v40n1Supl.65913}}
\end{softbox}

\begin{softbox}{milcTurquesa}{milcGris}{Saber contemporáneo}
\textbf{Monitorear} es medir la misma variable, con la misma frecuencia, durante un tiempo, y anotar cada valor. Un proyecto de monitoreo junta lo que hiciste en el periodo: leer sensores, calibrar umbrales, depurar y combinar señales. Tiene cuatro propiedades. (1) Una \textbf{variable bien definida}: la temperatura del aula, no «el clima». (2) Una \textbf{frecuencia constante}: cada hora, no cuando te acuerdes. (3) Un \textbf{registro que queda}: bitácora en papel o en una hoja de cálculo, con hora y observación. (4) Un \textbf{análisis al final}: releer las treinta filas buscando un patrón, una subida, un pico, una hora repetida. Con el micro:bit, el registro es sencillo: el programa muestra el valor y el nivel de alerta, y tú anotas a la hora acordada. Sin propósito, el monitoreo es acumular números. Con propósito, es buscar el dato que cambia una decisión.
\end{softbox}

\begin{softbox}{milcMostaza}{milcGris}{Pregunta-puente}
\textit{El caficultor solo ve el ciclo completo mirando el año, no una semana. Si midieras la temperatura del salón una sola vez a las 10, ¿qué te perderías de lo que pasa a la 1?}
\end{softbox}

\begin{softbox}{milcVerde}{milcGris}{Saber y saber hacer}
Al terminar podrás: (1) \textbf{identificar} cinco variables del colegio que valga la pena medir y la decisión que saldría de cada una; (2) \textbf{crear} el programa y el plan de tres jornadas con frecuencia constante; (3) \textbf{evaluar} la primera jornada de mediciones y ajustar lo que falló; (4) \textbf{explicar} por qué treinta datos en tres días valen más que cien en una hora.
\end{softbox}



\small
\begin{tabularx}{\textwidth}{>{\bfseries}p{3.0cm}Y}
\rowcolor{milcGradePrimary!12}Autor & Dr. Álvaro Cárdenas Orozco\\
DBA & Pensamiento computacional aplicado a sistemas de control con sensores y actuadores (MEN, T\&I)\\
Referentes & Pensamiento computacional · MakeCode / micro:bit · Logos como razón universal\\
Producto final & Un micro:bit programado para medir una variable real del colegio, temperatura del aula, luz del salón, ruido del recreo o paso de personas, durante tres jornadas. Una bitácora con al menos treinta mediciones, diez por jornada, con hora, valor, observación y nivel de alerta. Y un análisis de un patrón, la cifra que lo respalda y una propuesta de acción que salga de los datos.\\
\end{tabularx}
\normalsize

\puente{En la sesión 8 tu sistema decidió con lógica compuesta. Hoy lo dejas trabajando solo durante tres jornadas y anotas lo que ve. En la sesión 10 vas a sustentar el proyecto con los datos en la mano.}

\milcestacion{milcGradeDark}{Ruta MILC: así vamos a aprender}
\begin{tabularx}{\textwidth}{>{\centering\arraybackslash}X>{\centering\arraybackslash}X>{\centering\arraybackslash}X>{\centering\arraybackslash}X}
\phasecard{1}{milcVerde}{milcGris}{Escucha\\[-1mm]\small Observo, escucho y pregunto.} &
\phasecard{2}{milcTurquesa}{milcGris}{Entender\\[-1mm]\small Sistematizo y ordeno.} &
\phasecard{3}{milcMagenta}{milcGris}{Hacer\\[-1mm]\small Construyo y aplico.} &
\phasecard{4}{milcVino}{milcGris}{Cierre\\[-1mm]\small Evalúo y reflexiono.}
\end{tabularx}


\puente{Antes de programar, vas a listar cinco variables del colegio y preguntarte qué decisión saldría de cada una. Sin decisión no hay proyecto: hay una lista de números.}

\milcestacion{milcVerde}{Estación 1 de 3 · Escucha}

\begin{softbox}{milcVerde}{milcGris}{Escena para escuchar}
\textbf{Actividad 1 · IDENTIFICA --- Cinco variables y una decisión} (15 min · individual). (1) Anota cinco variables del colegio que un micro:bit pueda medir durante varios días: temperatura del salón, luz a distintas horas, ruido de la cafetería, paso de personas frente al baño, humedad de una matera del huerto. (2) Para cada una, escribe qué sensor usarías. (3) Para cada una, escribe qué decisión podría salir del dato: abrir ventanas a cierta hora, mover una clase, regar menos. (4) Escribe cada cuánto la medirías. (5) Elige una y escribe en una línea por qué esa.
\end{softbox}

{\sffamily\bfseries\fontsize{8}{10}\selectfont\textcolor{milcNegro!55}{MARCA CUANDO LO HAYAS HECHO}}
\milccheckrow{Tengo cinco variables con su sensor.}
\milccheckrow{Cada variable tiene una decisión que saldría del dato.}
\milccheckrow{Elegí una y escribí por qué.}

\begin{infoband}{milcVerde}


\textbf{Lo que va al cuaderno (Actividad 1 · IDENTIFICA).} \textbf{Título:} «Actividad 1 --- Cinco variables y una decisión». \textbf{Formato:} tabla de 5 filas y 4 columnas: variable, sensor, decisión posible, frecuencia; y la elegida con su porqué. \textbf{Extensión:} media página. \textbf{Sabes que terminaste cuando} la variable elegida tiene una decisión concreta que alguien del colegio podría tomar con el dato.
\end{infoband}

\puente{Ya elegiste la variable y la decisión. Ahora aprendes las cuatro propiedades del monitoreo y, con tu pareja, programas el micro:bit y armas el plan de tres jornadas.}

\milcestacion{milcTurquesa}{Estación 2 de 3 · Entender}

\begin{softbox}{milcTurquesa}{milcGris}{Lo que vas a entender}
\textbf{Actividad 2 · CREA --- El programa y el plan de tres jornadas} (30 min · parejas). (1) Con tu pareja, en MakeCode, guarden la lectura del sensor elegido en una variable y muéstrenla con «mostrar número». (2) Agreguen las tres ramas con los umbrales de la sesión 6 y un ícono por nivel, para que el micro:bit muestre valor y alerta. (3) Escriban el plan: dónde va a estar el micro:bit, a qué horas se anota, quién anota cada día. (4) Dibujen la bitácora de cuatro columnas: hora, valor, observación del entorno y nivel de alerta. (5) Fijen la frecuencia: diez mediciones por jornada, a horas que se puedan cumplir.
\end{softbox}

\milcpaso{1}{milcTurquesa}{\textbf{Una variable, una unidad.} Temperatura en grados Celsius; luz de 0 a 255; ruido en el nivel que da el micrófono; personas como conteo con «al agitar» o un botón. Si la unidad cambia entre filas, la bitácora no se puede comparar.}
\milcpaso{2}{milcTurquesa}{\textbf{Frecuencia que se cumpla.} Diez mediciones por jornada a horas fijas: al llegar, en cada cambio de clase, en el recreo, al salir. Mejor cada 45 minutos siempre, que cada 10 minutos el lunes y nada el martes.}
\milcpaso{3}{milcTurquesa}{\textbf{Tres jornadas, no una.} Día 1: instalar, medir y ajustar la calibración. Días 2 y 3: medir con la frecuencia fija. Un pico de calor a la 1 de la tarde solo es un patrón si aparece los tres días.}
\milcpaso{4}{milcTurquesa}{\textbf{Cómo se hace, paso a paso.} (1) Programar valor y alerta. (2) Escribir el plan con lugar, horas y responsable. (3) Dibujar la bitácora. (4) Medir tres jornadas. (5) Releer las treinta filas. (6) Escribir el patrón, la cifra y la propuesta.}

\begin{drawbox}{milcTurquesa}{El monitoreo, en una mirada}
\textbf{Variable:} una sola, con su unidad. \\[1mm] \textbf{Frecuencia:} cada cuánto, siempre igual. \\[1mm] \textbf{Bitácora:} hora, valor, observación, alerta. \\[1mm] \textbf{Patrón:} lo que se repite en las treinta filas. \\[1mm] \textbf{Propuesta:} la decisión que sale del patrón.
\end{drawbox}

\begin{softbox}{milcMostaza}{milcGris}{Errores comunes y su corrección}
(1) \textbf{Medir sin propósito}: treinta números y ninguna decisión. (2) \textbf{Frecuencia que cambia}: cada hora el lunes, dos veces el martes. (3) \textbf{Confiar en la memoria}: «después anoto» y la fila se pierde. (4) \textbf{No releer}: la bitácora llena y nadie la mira al final. (5) \textbf{Propuesta sin cifra}: «hace calor» en vez de «sube 3 grados entre 12 y 1».
\end{softbox}

\begin{infoband}{milcTurquesa}


\textbf{Lo que va al cuaderno (Actividad 2 · CREA).} \textbf{Título:} «Actividad 2 --- El programa y el plan de tres jornadas». \textbf{Formato:} el plan con lugar, horas y responsable por día, y la bitácora de cuatro columnas dibujada con treinta filas vacías numeradas. \textbf{Extensión:} una página. \textbf{Sabes que terminaste cuando} el micro:bit muestra valor y alerta, y cualquiera de los dos puede anotar una fila sin preguntarle al otro.
\end{infoband}

\puente{Ya tienes el programa y el plan. Toca la primera jornada: instalar, medir, anotar y ajustar lo que no funcionó. Las otras dos jornadas siguen esta semana.}

\milcestacion{milcMagenta}{Estación 3 de 3 · Hacer}

\begin{softbox}{milcMagenta}{milcGris}{Cómo vas a construir tu producto}
\textbf{Actividad 3 · EVALÚA --- La primera jornada} (25 min · individual). (1) Instala el micro:bit en el lugar del plan y toma la primera medición: hora, valor, observación, alerta. (2) Toma dos mediciones más en la sesión, con unos minutos entre ellas, cambiando algo del entorno: abre la ventana, apaga la luz. (3) Revisa si el nivel de alerta que salió tiene sentido con lo que ves; si no, ajusta el umbral y anótalo. (4) Escribe en una línea qué falló del plan: la hora, el lugar, la batería, el acceso. (5) Completa las diez mediciones de la jornada 1 en el resto del día y trae las jornadas 2 y 3 a la sesión 10.
\end{softbox}

\milcpaso{1}{milcMagenta}{\textbf{Tu entrega tiene tres piezas.} Hoy: el plan y las primeras filas de la jornada 1. En la sesión 10: la bitácora con treinta filas, el análisis con patrón y cifra, y la propuesta de acción. El programa se comparte con enlace o archivo .hex.}
\milcpaso{2}{milcMagenta}{\textbf{El análisis final tiene tres partes.} (1) El patrón en una frase: «la temperatura del salón sube entre las 12 y la 1 los tres días». (2) La cifra que lo respalda: «de 24 a 28 grados». (3) La propuesta: «abrir las ventanas a las 11:30». Sin cifra, el patrón es una impresión.}
\milcpaso{3}{milcMagenta}{\textbf{Lo que el ojo no ve.} Tú estás en el salón todos los días y no notas que sube 4 grados. La bitácora sí. Ese es el hallazgo que vas a contar: lo que treinta filas mostraron y la costumbre escondía.}
\milcpaso{4}{milcMagenta}{\textbf{Si algo falla, se anota.} Se acabó la batería, alguien movió el micro:bit, se perdieron dos horas. Eso va en la observación, no se disimula. Una bitácora con huecos declarados vale más que una completa inventada.}

\begin{drawbox}{milcMagenta}{Antes de entregar}
\checkbox Variable elegida con una decisión escrita. \\[1mm] \checkbox Micro:bit programado con valor y nivel de alerta. \\[1mm] \checkbox Plan con lugar, horas y responsable por jornada. \\[1mm] \checkbox Bitácora de cuatro columnas con las filas de la jornada 1. \\[1mm] \checkbox Umbral revisado en la primera jornada, con ajuste anotado si hizo falta. \\[1mm] \checkbox Una línea sobre qué falló del plan. \\[1mm] \checkbox Programa compartido con enlace o archivo .hex.
\end{drawbox}

\begin{softbox}{milcMagenta}{milcGris}{Plantilla de trabajo}
\textbf{Variable.} \textit{Mido … en …} \\[1mm] \textbf{Decisión.} \textit{Con el dato se podría …} \\[1mm] \textbf{Frecuencia.} \textit{Diez veces por jornada, a las …} \\[1mm] \textbf{Patrón.} \textit{Los tres días, … entre … y …} \\[1mm] \textbf{Cifra.} \textit{De … a …} \\[1mm] \textbf{Propuesta.} \textit{Por eso propongo …}
\end{softbox}

\begin{infoband}{milcMagenta}


\textbf{Lo que va al cuaderno (Actividad 3 · EVALÚA).} \textbf{Título:} «Actividad 3 --- La primera jornada». \textbf{Formato:} las tres primeras filas de la bitácora, el ajuste del umbral si lo hubo, y la línea sobre qué falló del plan. \textbf{Extensión:} media página. \textbf{Sabes que terminaste cuando} las tres filas tienen hora, valor, observación y alerta, y el nivel de alerta coincide con lo que viste en el salón.
\end{infoband}

\puente{Lo que entregas en la sesión 10: la bitácora de treinta mediciones, el patrón con su cifra y la propuesta. Lo que entregas hoy: el plan y las diez primeras filas. La prueba de calidad: cada fila tiene hora, valor y observación.}

\milcestacion{milcMostaza}{Producto de la sesión}
\textbf{Micro:bit midiendo tres jornadas + bitácora de treinta mediciones + patrón con cifra + propuesta de acción}.
\begin{softbox}{milcMostaza}{milcGris}{Criterios de calidad}
(1) La variable tiene una unidad clara y una decisión escrita antes de medir. (2) El micro:bit muestra el valor y el nivel de alerta con umbrales calibrados. (3) El plan fija lugar, horas y responsable, y la frecuencia se cumple las tres jornadas. (4) La bitácora tiene al menos treinta filas con hora, valor, observación y alerta. (5) El análisis dice un patrón en una frase y la cifra que lo respalda. (6) La propuesta de acción sale del patrón y alguien del colegio podría tomarla.
\end{softbox}

\puente{Antes de cerrar, mira lo que hiciste desde cinco lados. Un proyecto con propósito escrito se puede sustentar; una lista de números sin pregunta, no.}

\milcestacion{milcVino}{Evaluación liberadora · cinco dimensiones}

\begin{softbox}{milcVino}{milcGris}{Sentido de la evaluación}
Evaluar no es solo revisar una nota. Es reconocer qué aprendí, qué cambió en mí, qué emociones manejé, cómo me relaciono con la ciudad y la comunidad, y qué le dejaré a quien viene detrás.
\end{softbox}

\textbf{Mi reflexión liberadora en cinco dimensiones.} Completa cada frase iniciada:

\small
\begin{tabularx}{\textwidth}{>{\bfseries\RaggedRight\arraybackslash}p{3.4cm}Y>{\RaggedRight\arraybackslash}p{4.6cm}}
\rowcolor{milcVino}\color{white}Dimensión & \color{white}\bfseries Mi respuesta (frame starter) & \color{white}\bfseries Algo que descubrí\\
1. Personal & Lo que más cambió en mí fue \linead{6.5cm} & \linead{4.2cm}\\
2. Emocional & La emoción más fuerte fue \linead{2.5cm} y la regulé \linead{3cm} & \linead{4.2cm}\\
3. Ciudadana & En democracia esto importa porque \linead{6cm} & \linead{4.2cm}\\
4. Local (Cartago) & En mi barrio sería útil para \linead{6.5cm} & \linead{4.2cm}\\
5. Intergeneracional & A mi abuela/abuelo le diría \linead{6.5cm} & \linead{4.2cm}\\
\end{tabularx}
\normalsize

\begin{infoband}{milcVino}
\textbf{Para la próxima sustentación.} Vuelve a esta tabla en seis meses. Verás que las respuestas cambian — no porque lo aprendido cambie, sino porque tú cambias. La evaluación liberadora es una conversación contigo a través del tiempo.
\end{infoband}


\puente{Para terminar, tres ideas sobre el tiempo que hace falta para ver, contadas con palabras de esta guía: Dussel sobre perder tiempo para elegir lo importante, Marco Aurelio sobre mirar los ciclos, y Floridi sobre los patrones pequeños.}

\milcestacion{milcGradeDark}{Tres ideas para cerrar · Dussel, un estoico y Floridi}

\begin{softbox}{milcVino}{milcGris}{Enrique Dussel · lente del nosotros}
\textit{Como los temas son infinitos y el tiempo corto, hay que saber perder tiempo para elegir lo que de verdad importa.}

{\footnotesize\itshape\color{milcNegro!70}--- idea tomada de Enrique Dussel, contada con palabras de esta guía. No es una cita textual.}

Podrías medir cinco variables a la vez y no entender ninguna. Elegir una, con una decisión escrita, es «perder» quince minutos para que las tres jornadas sirvan de algo.

\textbf{¿Qué decisión de mi colegio cambiaría si alguien tuviera estos treinta datos?}
\end{softbox}
\linead{\linewidth}\\[1mm]
\linead{\linewidth}

\begin{softbox}{milcMostaza}{milcGris}{Estoicismo · Marco Aurelio · cuidado interior}
\textit{Conviene mirar el curso de los astros y las transformaciones de las cosas, como quien gira con ellas.}

{\footnotesize\itshape\color{milcNegro!70}--- idea tomada de Marco Aurelio, contada con palabras de esta guía. No es una cita textual.}

El caficultor mira el año para ver las dos cosechas y el tiempo frío. Tú miras tres jornadas para ver el pico de la 1 de la tarde. Lo que se repite solo aparece cuando miras el ciclo entero.

\textbf{¿Qué de mi día se repite sin que lo note, porque nunca lo miré tres veces seguidas?}
\end{softbox}
\linead{\linewidth}\\[1mm]
\linead{\linewidth}

\begin{softbox}{milcTurquesa}{milcGris}{Luciano Floridi · lente de la infoesfera}
\textit{Los patrones pequeños solo significan algo si se juntan bien, se comparan y se procesan a tiempo.}

{\footnotesize\itshape\color{milcNegro!70}--- idea tomada de Luciano Floridi, contada con palabras de esta guía. No es una cita textual.}

Una medición de 28 grados no dice nada. Treinta mediciones, comparadas por hora y por día, dicen «sube a la 1». Y solo sirven si las relees antes de la sesión 10, no después.

\textbf{¿Qué patrón pequeño de mi bitácora se perdería si no comparo los tres días?}
\end{softbox}
\linead{\linewidth}\\[1mm]
\linead{\linewidth}

\begin{infoband}{milcNegro}
\textbf{Nota para el docente.} Las tres ideas están contadas con palabras de esta guía a partir del pensamiento de cada autor; no son citas textuales. De 9.º en adelante se trabajan con la obra y su referencia.
\end{infoband}

\milcestacion{milcVino}{Mi autoevaluación · marca una casilla por fila}

{\small
\setlength{\extrarowheight}{2.2pt}
\begin{tabularx}{\textwidth}{>{\RaggedRight\arraybackslash\bfseries}p{3.0cm}YYY}
\rowcolor{milcVino}\color{white}\bfseries Criterio & \color{white}\bfseries Lo logré & \color{white}\bfseries En proceso & \color{white}\bfseries Necesito apoyo\\[.6mm]
Comprensión del tema & \milcopcion{Explico las ideas con mis palabras.} & \milcopcion{Reconozco las ideas, falta precisión.} & \milcopcion{Necesito repasar lo básico.}\\[.8mm]
\rowcolor{milcGris}Pensamiento computacional & \milcopcion{Usé los pilares al armar mi producto.} & \milcopcion{Usé algunos pilares.} & \milcopcion{Necesito releer la estación 2.}\\[.8mm]
Aplicación y producto & \milcopcion{Mi producto cumple los criterios.} & \milcopcion{Está armado, con ajustes pendientes.} & \milcopcion{Aún está en construcción.}\\[.8mm]
\rowcolor{milcGris}Reflexión 5 dimensiones & \milcopcion{Respondí cada una con honestidad.} & \milcopcion{Respondí algunas con detalle.} & \milcopcion{Mis respuestas son muy generales.}\\[.8mm]
Triángulo de pensamiento & \milcopcion{Conecté las 3 voces con mi trabajo.} & \milcopcion{Conecté al menos una.} & \milcopcion{No las relaciono con mi proceso.}\\[.8mm]
\end{tabularx}}

\vspace{3mm}
\textbf{Hoy entendí que:}\\[1mm]
\linead{\linewidth}\\[1mm]
\linead{\linewidth}

\textbf{Mi compromiso para la próxima sesión es:}\\[1mm]
\linead{\linewidth}\\[1mm]
\linead{\linewidth}

\begin{softbox}{milcMostaza}{milcGris}{Cita de despedida · Marco Aurelio}
\textit{``El obstáculo en el camino se vuelve el camino. Lo que estorba al camino se vuelve camino.''}\\[1mm]
Cada error de hoy, bien observado, se convierte en pieza del camino que pisarás mañana.
\end{softbox}

\small
\begin{tabularx}{\textwidth}{>{\bfseries}p{3.6cm}Y}
\rowcolor{milcGradeDark!12}Nombre completo: & \linead{10cm}\\
Fecha: & \linead{4cm} \quad Curso/grupo: \linead{4cm}\\
Mi firma: & \linead{9cm}\\
\end{tabularx}
\normalsize

\begin{infoband}{milcGradeDark}
\textbf{Para la próxima sesión.} Completa las jornadas 2 y 3 y llega con la bitácora de treinta filas, el patrón con su cifra y la propuesta. En la sesión 10 vas a sustentar el proyecto en cinco minutos con el micro:bit funcionando. Los datos que anotes esta semana son lo que vas a defender.
\end{infoband}

\end{document}
